\Article{25}

\textsc{Theorema.}\quad \textit{Sit $n$ productum e duobus integris positivis inter se primis $a$, $b$,} \textit{statuaturque}

\[
\begin{aligned}
P&=1+ r^{aa}+r^{4aa}+r^{9aa}+\text{etc.}+r^{(b-1)^2aa}\\
Q&=1+r^{bb}+r^{4bb}+r^{9bb}+\text{etc.}+r^{(a-1)^2bb}
\end{aligned}
\]

\textit{Tum dico fore $W=PQ$.}

\textit{Demonstr.}\quad Designet $\alpha$ indefinite numeros $0,1,2,3\dotsfour a-1$, $\beta$ indefinite numeros $0,1,2,3\dotsfour b-1$, $\nu$ indefinite numeros $0,1,2,3\dotsfour n-1$.\quad Tunc patet esse

\[
\begin{aligned}
P=\Sigma r^{aa\beta\beta},\quad Q=\Sigma r^{bb\alpha\alpha},\quad W=\Sigma r^{\nu\nu}
\end{aligned}
\]

Hinc erit $PQ=\Sigma r^{aa\beta\beta+bb\alpha\alpha}$, substituendo pro $\alpha$ et $\beta$ omnes valores, omnibus modis inter se combinatos; hinc porro propter $2ab\alpha\beta=2\alpha\beta n$, erit $PQ=\Sigma r^{(a\beta+b\alpha)^2}$. Sed nullo negotio perspicitur, singulos valores ipsius $a\beta+b\alpha$ inter se diversos esse, atque alicui valori ipsius $\nu$ aequales.\quad Hinc erit $PQ=\Sigma r^{\nu\nu}=W$.

Ceterum notandum est, $r^{aa}$ esse radicem propriam aequationis $x^b-1=0$, atque $r^{bb}$ radicem propriam aequationis $x^a-1=0$.

\SourceLeaf{47}

\Article{26}

Sit porro $n$ productum e tribus numeris inter se primis $a,b,c$, patetque, si statuatur $bc=b'$, etiam $a$ et $b'$ inter se primos fore; adeoque $W$ productum e duobus factoribus

\[
\begin{aligned}
&1+r^{aa}+r^{4aa}+r^{9aa}+\text{etc.}+r^{(b'-1)^2aa}\\
&1+r^{b'b'}+r^{4b'b'}+r^{9b'b'}+\text{etc.}+r^{(a-1)^2b'b'}
\end{aligned}
\]

Sed quum $r^{aa}$ sit radix propria aequationis $x^{bc}-1=0$, erit ipse factor prior productum ex

\[
\begin{aligned}
&1+\rho^{bb}+\rho^{4bb}+\rho^{9bb}+\text{etc.}+\rho^{(c-1)^2bb}\\
&1+\rho^{cc}+\rho^{4cc}+\rho^{9cc}+\text{etc.}+\rho^{(b-1)^2cc}
\end{aligned}
\]

si statuitur $r^{aa}=\rho$.\quad Hinc patet, $W$ esse productum e factoribus tribus

\[
\begin{aligned}
&1+r^{bbcc}+r^{4bbcc}+r^{9bbcc}+\text{etc.}+r^{(a-1)^2bbcc}\\
&1+r^{aacc}+r^{4aacc}+r^{9aacc}+\text{etc.}+r^{(b-1)^2aacc}\\
&1+r^{aabb}+r^{4aabb}+r^{9aabb}+\text{etc.}+r^{(c-1)^2aabb}
\end{aligned}
\]

ubi $r^{bbcc}$, $r^{aacc}$, $r^{aabb}$ erunt resp. radices propriae aequationum $x^a-1=0$, $x^b-1=0$, $x^c-1=0$.

\Article{27}

Hinc facile concluditur generaliter, si $n$ sit productum e factoribus quotcunque inter se primis $a,b,c$ etc., $W$ fieri productum e totidem factoribus, qui sint

\[
\begin{aligned}
&1+r^{\frac{nn}{aa}}+r^{\frac{4nn}{aa}}+r^{\frac{9nn}{aa}}+\text{etc.}+r^{\frac{(a-1)^2nn}{aa}}\\
&1+r^{\frac{nn}{bb}}+r^{\frac{4nn}{bb}}+r^{\frac{9nn}{bb}}+\text{etc.}+r^{\frac{(b-1)^2nn}{bb}}\\
&1+r^{\frac{nn}{cc}}+r^{\frac{4nn}{cc}}+r^{\frac{9nn}{cc}}+\text{etc.}+r^{\frac{(c-1)^2nn}{cc}},\;\text{etc.}
\end{aligned}
\]

ubi $r^{\frac{nn}{aa}}$, $r^{\frac{nn}{bb}}$, $r^{\frac{nn}{cc}}$ etc. erunt radices propriae aequationum $x^a-1=0$, $x^b-1=0$. $x^c-1=0$ etc.

\Article{28}

Ex his principiis transitus ad determinationem completam ipsius $W$ pro valore quocunque ipsius $n$ sponte iam obvius est.\quad Decomponatur scilicet $n$ in facto\SourceLeaf{48}

res $a,b,c$ etc. tales, qui sint vel numeri primi inaequales, vel potestates numerorum primorum inaequalium, statuatur $r^{\frac{nn}{aa}}=A$, $r^{\frac{nn}{bb}}=B$, $r^{\frac{nn}{cc}}=C$ etc. eruntque $A,B,C$ etc. radices propriae aequationum $x^a-1=0$, $x^b-1=0$, $x^c-1=0$ etc., atque $W$ productum e factoribus

\[
\begin{aligned}
&1+A+A^4+A^9+\text{etc.}+A^{(a-1)^2}\\
&1+B+B^4+B^9+\text{etc.}+B^{(b-1)^2}\\
&1+C+C^4+C^9+\text{etc.}+C^{(c-1)^2}
\end{aligned}
\]

Sed hi singuli factores per ea, quae in artt. $20,21,23$ docuimus, determinari poterunt, unde etiam valor producti innotescet.\quad Regulas pro determinandis illis factoribus hic in unum obtutum collegisse haud inutile erit.\quad Quum radix $A$ fiat $=\frac{kn}{a}\pdot\frac{360^0}{a}$, aggregatum $1+A+A^4+A^9+\text{etc.}+A^{(a-1)^2}$, quod per $L$ denotabimus, perinde per numerum $\frac{kn}{a}$ determinabitur, ut in disquisitione nostra generali $W$ per $k$.\quad Duodecim iam casus sunt distinguendi.

I.\quad Si $a$ est numerus primus formae $4\mu+1$, puta $=p$, vel potestas talis numeri primi cum exponente impari, simulque $\frac{kn}{a}$ residuum quadraticum ipsius $p$, erit $L=+\surd a$.

II.\quad Si manentibus reliquis $\frac{kn}{a}$ est non-residuum quadraticum ipsius $p$, erit $L=-\surd a$.

III.\quad Si $a$ est numerus primus formae $4\mu+3$, puta $=p$, vel potestas talis numeri primi cum exponente impari, simulque $\frac{kn}{a}$ residuum quadraticum ipsius $p$, erit $L=+i\surd a$.

IV.\quad Si, manentibus reliquis ut in III, $\frac{kn}{a}$ est non-residuum quadraticum ipsius $p$, erit $L=-i\surd a$.

V.\quad Si $a$ est quadratum, altiorve potestas numeri primi (imparis) cum exponente pari, erit $L=+\surd a$.

VI.\quad Si $a=2$, erit $L=0$.

VII.\quad Si $a=4$, altiorve potestas binarii cum exponente pari, simulque $\frac{kn}{a}$ formae $4\mu+1$, erit $L=(1+i)\surd a$.

VIII.\quad Si, manentibus reliquis ut in VII, $\frac{kn}{a}$ est formae $4\mu+3$, erit $L=(1-i)\surd a$.

IX.\quad Si $a=8$, altiorve potestas binarii cum exponente impari, simulque $\frac{kn}{a}$ formae $8\mu+1$, erit $L=(1+i)\surd a$.

\SourceLeaf{49}

X.\quad Si, manentibus reliquis ut in IX, $\frac{kn}{a}$ est formae $8\mu+3$, erit $L=(-1+i)\surd a$.

XI.\quad Si manentibus reliquis $\frac{kn}{a}$ est formae $8\mu+5$, erit $L=(-1-i)\surd a$.

XII.\quad Si manentibus reliquis $\frac{kn}{a}$ est formae $8\mu+7$, erit $L=(1-i)\surd a$.

\Article{29}

Sit exempli caussa $n=2520=8\pdot9\pdot5\pdot7$, atque $k=13$.\quad Hic erit

{\small
\[
\begin{aligned}
&\text{pro }a=8,\;\text{per casum XII,}\quad L=(1-i)\surd8\\
&\text{pro factore }9,\;\text{per casum V, summa respondens erit }=\surd9\\
&\text{pro factore }5,\;\text{per casum II, summa respondens erit }=-\surd5\\
&\text{pro factore }7,\;\text{per casum III, summa respondens erit }=+i\surd7\mathord{.}
\end{aligned}
\]
}

Hinc fit $W=(1-i)\pdot(-i)\pdot\surd2520=(-1-i)\surd2520$.

Sit pro eodem valore ipsius $n$, $k=1$: tunc respondebit

\[
\begin{aligned}
&\text{factori }8\;\text{summa }(-1+i)\surd8\\
&\text{factori }9\;\text{summa }\surd9\\
&\text{factori }5\;\text{summa }\surd5\\
&\text{factori }7\;\text{summa }-i\surd7
\end{aligned}
\]

Hinc conflatur productum $W=(1+i)\surd2520$.
